%!TEX root = ../Thesis.tex
\chapter{An Appendix}

\subsection{Examples of hyperparameter configurations}
\label{appendix:hyperparam}

\begin{figure}[htp]
\centering
\figcross
\begin{enumerate}
    \item $C$: the constant subtracted from the Gaussian-weighted sum to compute the local threshold (see \cref{subsubsec:gaussian_baseline}),
    \item $d_{gaussian}$: the diameter of the kernel window for Gaussian adaptive thresholding (see \cref{subsubsec:gaussian_baseline}),
    \item $d_{bilateral}$: the diameter of the kernel window for bilateral pre-filtering (see \cref{subsubsec:gaussian_baseline}),
    \item $\sigma$: the spatial and intensity standard deviation for bilateral pre-filtering (see \cref{subsubsec:gaussian_baseline}),
\end{enumerate}
\caption{Hyperparameters for \acf{Gaussian}.}
\label{fig:hyperparameters_gaussian}
\end{figure}

\begin{figure}[htp]
\centering
\figcross
\begin{enumerate}
    \item $\sigma$: the constant used to determine the thresholds $T_l$ and $T_h$ for the hysteresis thresholding,
    \item $d_{closing}$: the diameter of the kernel window for the morphological closing of the edges
\end{enumerate}
\caption{Hyperparameters for \acf{CannyFill}.}
\label{fig:hyperparameters_CannyFill}
\end{figure}

\begin{figure}
\centering
\figcross
\begin{enumerate}
    \item $C$: the constant subtracted from the Gaussian-weighted sum to compute the local threshold (see \cref{subsubsec:gaussian_baseline}),
    \item $d_{gaussian}$: the diameter of the kernel window for Gaussian adaptive thresholding (see \cref{subsubsec:gaussian_baseline}),
    \item $d_{bilateral}$: the diameter of the kernel window for bilateral pre-filtering (see \cref{subsubsec:gaussian_baseline}),
    \item $\sigma$: the spatial and intensity standard deviation for bilateral pre-filtering (see \cref{subsubsec:gaussian_baseline}),
    \item $n_{fgd\_points}$: the number of foreground point prompts to be sampled from the foreground of the first Gaussian thresholded mask,
    \item $n_{bgd\_points}$: the number of background point prompts to be sampled from the sure background region,
    \item $d_{gap\_erosion}$: the kernel size for erosion used to create a gap between the foreground and background regions.
\end{enumerate}
\caption{Hyperparameters for \acf{SAM+pts}}
\label{fig:hyperparameters_SAM}
\end{figure}

\clearpage


\section{Interactive tool development}

\subsection{Usability tests}
\label{appendix:usability-tests}

\subsubsection{Requirement analysis and first usability test}
The 5th of May, a requirement analysis and first usability test of a MVP was performed with Ester through a video call.
That MVP was composed of two tools: a classical method (Gaussian), and an interactive SAM version (MobileSAM).
After trying out the tools, she gave the impression that she could probably use the outputs as draft drawings to be completed with a pen, which might save her some time.

For the classical baseline, the grainy resolution (due to the input image quality) bothered her a bit.
The output mask was also noisy at times. And not all lines were always picked up.

For the SAM version, she pointed out that it did not pick up the noisy tablet damage as much as the classical method did.
But she was worried that the interaction with SAM would take longer than simply doing strokes manually.
The box-prompt workflow was not completely successful during the test.
Due to bugs, boxes mostly did not work. and when they did, the segmented forms were not useful.

\subsubsection{Second (small) usability test}

The 8th of May, another usability test was performed on an updated MVP with Ester through a video call.
Now, the box prompts worked, and the MobileSAM was replaced with SAM2.1.
The segmentation quality per box-prompted sign was noticeably better than previously.
Together with Lukas Esterle, the conclusions were to make a slider for the classical model to decrease and increase granularity.
It was also suggested that box prompts for SAM could be computationally preselected based on the mask from the classical model, but this was discarded as it was not deemed needed.

\subsubsection{Third usability and final acceptance test}

The 12th of May, a third usability test was conducted in person at AIAS in Århus, including the first part of the usability experiment.
As all 45 images could not be drawn all at once, the .ora exports from the tool were saved, and finished at a later time in batches by Ester on her own, for each of them recording the time it took to do so.
Due to the huge amount of time it took to process the images with the tool (especially with the SAM version), the experiment was stopped after just 3 images.
